\section{模型假设与基本约定}

为了简化问题，便于模型的建立和求解，我们作出如下合理假设：

\begin{enumerate}
  \item 假设各硬模块均为尺寸固定的刚性模块，其尺寸由 \texttt{.blocks} 文件中给出的顶点坐标确定，在布局过程中不发生缩放、剪切或形变。
  \item 假设问题一至问题三中的矩形模块只允许保持原方向或旋转 $90^\circ$，旋转不改变模块的实际面积；对于正方形模块，两种旋转状态视为同一种摆放方式。
  \item 假设不同模块的边界可以相互接触，但其内部区域不能发生重叠。由于整体平移不会改变布局面积、长宽比和模块间的相对关系，可以将芯片轮廓的左下角统一置于坐标原点。
  \item 假设问题一仅考虑硬模块的几何布局，不考虑模块之间的网络连接关系，因此 \texttt{.pl} 文件中的 Terminal 不参与模块位置、旋转及轮廓面积的计算。
  \item 假设问题二和问题三中各硬模块的网络引脚位于模块几何中心，外部 Terminal 的位置由 \texttt{.pl} 文件给定且保持不变。Terminal 仅参与总半周长线长 HPWL 的计算，不参与模块面积及模块间非重叠约束。
  \item 假设问题二中模块的放置坐标取整数，固定正方形轮廓由题目给定的死区率和模块总面积确定；问题三进一步将正方形轮廓边长限制为正整数，并以硬模块总面积为基础计算对应的死区率。
  \item 假设问题四中的 L 型、T 型等非矩形模块均由若干轴对齐的规则区域组成，其真实占用区域可以用单位格集合或有限个矩形区域的并准确表示；模块允许进行 $0^\circ$、$90^\circ$、$180^\circ$ 和 $270^\circ$ 旋转，旋转后几何形态相同的状态视为同一种形态。
  \item 对问题四题图中的尺寸标注按照统一比例进行解释，各模块的形状和尺寸在布局过程中保持不变；若题目对非矩形模块尺寸的定义发生改变，则相应布局结果需重新计算。
\end{enumerate}
